\section{Subestação}
-

\subsection{Dimensões}

\subsection{Altura da Subestação}

A altura da subestação é obtida através da somatório abaixo:

\begin{equation}
H_{se} = H_t + H_{ac} + H_c + H_i + H_{ab}
\end{equation}

$H_{se}$ - altura da subestação.
$H_t$ - altura do total do transformador, disponível na folha de especificação.
$H_{ac}$ - afastamento da chavee seccionadora.
$H_c$ - altura da chave seccionadora.
$H_i$ - altura do isolador.
$H_{ab}$ - afastamento do barramento

Usando os valores fornecidos pelo fabricantes, obtemos:

\begin{equation}
H_{se} = 1855 + 300 + 600 + 250 + 100 = 3105
\end{equation}

No entanto, por ser tratar de uma subestação com entrada aérea, a mesma deverá ter uma altura mínima de 5,5 m.

No projeto será construída uma subestação com altura de 6 metros.

\subsection{Paredes Externas}

Por ter, altura igual ou maior a 6 metros, a subestação terá paredes externas de 30cm e paredes de divisões internas de 15cm em alvenaria.


\subsection{Posto de Medição}

O posto de medição ocupará o espaço de 1,6m x 2,4m.

\subsection{Posto de Proteção}

A dimensão do posto de medição foi obtida por:

\begin{equation}
D_{cp} = D_{d} + 1000mm = 700 + 1000 = 1700mm
\end{equation}

\subsection{Posto de Transformação}

A dimensão do posto de transformação foi obtida por:

\begin{equation}
D_{ct} = D_{t} + 1000mm = 1985 + 1000 = 2985mm
\end{equation}

\subsection{Porta de acesso principal}

A norma especifica 1,20m x 2,10m como dimensão mínima para a porta de acesso pricipal. No projeto será usada portas de 1,40m x 2,10m.	

\subsection{Barramento Primário}

Será utilizada Barra de Cobre, 3/4" x 3/16", como especificado em norma.

